%to have line numbers
%\RequirePackage{lineno}
\documentclass[10pt, letterpaper]{article}      
\usepackage[margin=.1cm,font=small,labelfont=bf]{caption}[2007/03/09]
%\usepackage{endnotes}
%\let\footnote=\endnote
\usepackage{setspace}
\usepackage{longtable}                        
\usepackage{anysize}                          
\usepackage{natbib}                           
%\bibpunct{(}{)}{,}{a}{,}{,}                   
\bibpunct{(}{)}{,}{a}{}{,}                   
\usepackage{amsmath}
\usepackage[% draft,
pdftex]{graphicx} %draft is a way to exclude figures                
\usepackage{epstopdf}
\usepackage{hyperref}                             % For creating hyperlinks in cross references
\hypersetup{pdfborder={0 0 0.4}} %have nice light boxes on refs
 
% \usepackage[margins]{trackchanges}

% \note[editor]{The note}
% \annote[editor]{Text to annotate}{The note}
%    \add[editor]{Text to add}
% \remove[editor]{Text to remove}
% \change[editor]{Text to remove}{Text to add}

%TODO make it more standard before submission: \marginsize{2cm}{2cm}{1cm}{1cm}
\marginsize{1cm}{1cm}{.5cm}{.5cm}%{left}{right}{top}{bottom}   
					          % Helps LaTeX put figures where YOU want
 \renewcommand{\topfraction}{1}	                  % 90% of page top can be a float
 \renewcommand{\bottomfraction}{1}	          % 90% of page bottom can be a float
 \renewcommand{\textfraction}{0.0}	          % only 10% of page must to be text

 \usepackage{float}                               %latex will not complain to include float after float

\usepackage[table]{xcolor}                        %for table shading
\definecolor{gray90}{gray}{0.90}
\definecolor{orange}{RGB}{255,128,0}

\renewcommand\arraystretch{.9}                    %for spacing of arrays like tabular

%-------------------- my commands -----------------------------------------
\newenvironment{ig}[1]{
\begin{center}
 %\includegraphics[height=5.0in]{#1} 
 \includegraphics[height=3.3in]{#1} 
\end{center}}

 \newcommand{\cc}[1]{
\hspace{-.13in}$\bullet$\marginpar{\begin{spacing}{.6}\begin{footnotesize}\color{blue}{#1}\end{footnotesize}\end{spacing}}
\hspace{-.13in} }

%-------------------- END my commands -----------------------------------------



%-------------------- extra options -----------------------------------------

%%%%%%%%%%%%%
% footnotes %
%%%%%%%%%%%%%

%\long\def\symbolfootnote[#1]#2{\begingroup% %these can be used to make footnote  nonnumeric asterick, dagger etc
%\def\thefootnote{\fnsymbol{footnote}}\footnote[#1]{#2}\endgroup}	%see: http://help-csli.stanford.edu/tex/latex-footnotes.shtml

%%%%%%%%%%%
% spacing %
%%%%%%%%%%%

% \abovecaptionskip: space above caption
% \belowcaptionskip: space below caption
%\oddsidemargin 0cm
%\evensidemargin 0cm

%%%%%%%%%
% style %
%%%%%%%%%

%\pagestyle{myheadings}         % Option to put page headers
                               % Needed \documentclass[a4paper,twoside]{article}
%\markboth{{\small\it Politics and Life Satisfaction }}
%{{\small\it Adam Okulicz-Kozaryn} }

%\headsep 1.5cm
% \pagestyle{empty}			% no page numbers
% \parindent  15.mm			% indent paragraph by this much
% \parskip     2.mm			% space between paragraphs
% \mathindent 20.mm			% indent math equations by this much

%%%%%%%%%%%%%%%%%%
% extra packages %
%%%%%%%%%%%%%%%%%%

\usepackage{datetime}


\usepackage[latin1]{inputenc}
\usepackage{tikz}
\usetikzlibrary{shapes,arrows,backgrounds}


%\usepackage{color}					% For creating coloured text and background
%\usepackage{float}
\usepackage{subfig}                                     % for combined figures

\renewcommand{\ss}[1]{{\colorbox{blue}{\bf \color{white}{#1}}}}
\newcommand{\ee}[1]{\endnote{\vspace{-.10in}\begin{spacing}{1.0}{\normalsize #1}\end{spacing}\vspace{.20in}}}
\newcommand{\emd}[1]{\ExecuteMetaData[/tmp/tex]{#1}} % grab numbers  from stata

%TODO before submitting comment this out to get 'regular fornt'
\usepackage{sectsty}
\allsectionsfont{\normalfont\sffamily}
\usepackage{sectsty}
\allsectionsfont{\normalfont\sffamily}
\renewcommand\familydefault{\sfdefault}

%\usepackage[margins]{trackchanges}
\usepackage{rotating}
\usepackage{catchfilebetweentags}

\usepackage{abstract}
\renewcommand{\abstractname}{}    % clear the title
\renewcommand{\absnamepos}{empty} % originally center
%-------------------- END extra options -----------------------------------------
\date{{}\today \hspace{.2in}\xxivtime}
\title{  
  Personal Agency and Political Efficacy Across the Urban-Rural Gradient  
}
\author{
}

\begin{document}


\bibliographystyle{/home/aok/papers/root/tex/ecta}
\maketitle
\vspace{-.4in}
\begin{center}

\end{center}

%\noindent\textbf{target jou:}\\
 % cities if8.2,\\
  %ariq if3ish,\\
  %some soc j have soc lit: Toennies Simmel Durkheim maybe sth like j or urb soc or eur soc rev?\\
  %rubia: we'll need to add more lit for soc journals, maybe try cities first and then move down the list? we could also try urban studies? then if no luck we can add more lit and try soc journals.

\begin{abstract}
  \noindent This study uses World Values Survey 7 (2017-2023) across 47-62 countries
  to test urban-rural differences in personal agency and political efficacy.  
%
  Our findings indicate that personal agency is generally no higher in urban than in rural areas, contrary to the expectation of the ``city air makes you free'' hypothesis--there is mostly no difference. However, political efficacy is significantly higher in urban areas compared to rural areas in rich countries, consistent with the ``left behind ruralities'' hypothesis. As Toennies hypothesized, those doing well financially are ``active and alive'' in cities or, as we find here, have greater personal agency (but not political efficacy). Effect sizes are considerable at about ??. As in any correlational study, these findings do not establish causality. 
\end{abstract}
\vspace{.15in} 
\noindent{\sc World Values Survey, freedom, agency, autonomy, political efficacy, urbanism, urban, rural,
  urban-rural, cities
}
\vspace{.25in} 

\addtocontents{toc}{\protect\setcounter{tocdepth}{-2}}


\begin{spacing}{1.4} 
\rowcolors{1}{white}{gray90}


\noindent ``City life and Gesellschaft down the common people to decay and death.'' Ferdinand Toennies\\
%\citep[p. 231]{tonnies57}

\noindent ``In the city and therefore, where general conditions characteristic of the Gesellschaft prevail, only the upper strata, the rich and the cultured, are really active and alive.'' Ferdinand Toennies\\ %\citep[p. 227]{tonnies57}.

\section{Introduction}

``City air makes you free'' expresses a longstanding belief that cities offer freedom from the social constraints of rural life \citep{yamagishi12}. Hence, city residents are expected to be `freer' than those in rural areas. 
The saying referred originally to freedom from feudal obligations\citep{harvey12,park84}.
 \citep{yamagishi12}.\footnote{``Stadt Luft macht frei'' \citep[p. 12]{park84}. This saying originated in the Middle Ages, and it meant freedom from feudalism in cities; cities were non-feudal islands in a sea of feudalism \citep{harvey12}. Feudalism has long been gone, but we may be living under feudalism again, in the so called techno-feudalism \citep{varoufakis24,morozov2022critique}.
{Cloud  fiefdoms, rent, and digital serfdom through platforms like Google and Amazon ``enslave'' both urbanites and ruralities, but it's still perhaps possible that cities may be freer in the sense that there is more choice and opportunity with local traditional on the ground commerce, for example.}
%
While modern cities aren't clearly more free from techno-feudalism, as Middle Ages cities used to be from feudalism, cities are still often seen in the literature and popular culture as being more free than rural areas.}
Its contemporary appeal rests on the idea that cities provide more opportunities and choices associated with different life domains (i.e., how and where to work/study/shop/worship/live, whom to form relationships with). The question is whether these opportunities and choices also translate into greater perceived freedom and control over one's life. Finding little or no difference in personal agency between urban and rural residents would therefore challenge this expectation and be worth reporting.

To examine this question, we focus on personal agency, which refers to perceived freedom of choice and control over one's life. Many studies approach freedom as a predictor of subjective wellbeing, but freedom
is also an end-in-itself \citep{aok_free_from_to}, as recognized by 
\citet{fromm13,fromm12,fromm92,fromm64,fromm44,fromm94} and \citet{sen00}. It is
a desirable outcome in just about any respect, except perhaps for the
responsibility that comes with it. In the Middle Ages there was little freedom,
but also little responsibility--for instance, the poor and slaves were not
considered responsible for their lot--their plight was God-given. These days,
the poor worry about being poor and about it being their fault--everyone is
free, after all, and their experiences and achievements, or lack thereof, are
seen as being of their own making.\footnote{
 For instance,
 people in the U.S. feel relatively free and so they blame the poor for their
 lot, while people in Poland feel less free \citep{ee_ls15} and also blame the poor less
 (anecdotal evidence).} 
%in addition to
%free var could add poor are lazy etc from wvs  

Yet, feeling that one has control over one's own life does not necessarily mean feeling that one has a say in government decisions affecting one's community. This distinction bring us to our second variable, political efficacy.
%rubia: we need to add connection to political efficacy, otherwise feels disconnected.
Political efficacy refers to the belief that one has a say in government
decisions. The ``left behind ruralities'' hypothesis suggests that people in
rural areas may feel neglected by institutions and excluded from decisions
affecting their communities and day to day
life\citep{ulrich2021people,rodriguez2018revenge}. This suggests that rural
areas residents may also feel that they have less say in government decisions
than residents in cities. This may explain the growing rural discontent in
developed countries like the United States \citep{aok-swbGenYcity18}. 
Taken together, the ``city air makes you free'' and the ``left behind ruralities'' hypotheses raise the questions: do people in cities feel both greater control over their own lives and a greater say in government decisions, or do urban-rural differences depend on which aspect we consider?

In what follows, we attempt to answer these questions. First, we discuss the theoretical expectations, describe the measures and analytical approach--we measure freedom as (perceived) personal agency and then we turn to our other main variable of interest, (perceived) political efficacy--present our results, and consider their implications and limitations.  

\section{Literature}
\subsection{Perceived Personal Agency %Freedom/autonomy defined
}
 
% Some authors use term ``freedom'' \citep{aok_free_from_to} and some ``autonomy''
% \citep{steckermeier22}--we use the terms interchangeably, but prefer
% ``freedom'' as it seems to sound and fit better in this context.
 
What is ``perceived personal agency'' or ``freedom'' for short?\footnote{Here we focus on personal as opposed to
  societal freedom and follow \citet{aok_free_from_to} in this paragraph.} To be free is to be able to choose, to control and
direct one's own life. Freedom is having control--a person is free if she is her own master. It could be labeled `psychological freedom': ``Freedom means a degree of harmony between basic motives and overt behavior'' \citep[Bay 1970, p.83 cited in][]{aok_free_from_to}.
 It is the ``capability to choose,'' having ``awareness of alternatives'' and ``courage to choose'' \citep{veenhoven00_b}.\footnote{The degree of perceived control over choice is called ``locus of control'' in psychology and \citet{verme09} argues that the WVS item used here measures locus of control in addition to freedom.}
 Cities offer more choices than rural areas and, in this sense, more freedom. Yet, it is important to underscore that too much choice may also backfire and make it difficult to act or to choose, possibly leading to paralysis \citep{schwartz04}.

Freedom is a broad concept and could be sub-divided into dimensions--\citet{hitlin09} review the concept in sociology (refer also to a recent scale \citep{agencyScale25}). \citet{kouba16} provide another useful review, which is relevant here as it points to our WVS measure. Beyond its different dimensions, freedom also has value in its own right. The Universal Declaration of Human Rights recognizes freedom as a human right essential for the realization of one's human potential (art. 26, 29).

There are theoretical reasons to expect more freedom in cities than in rural areas. In rurality, there is mechanical solidarity--rural residents function like parts of a simple machine--groups are held together through homogeneity and similarity:
similar lifestyles, roles, shared values, and morality \citep{durkheim97}. Hence, rural areas enforce social conformity that is not conducive to individual freedom. On the other hand, organic solidarity in cities allows for more differentiation and individualism. %and individuation. %(Jung)
 Urbanites are also more tolerant \citep{tuch87}.
 % 
% 
Simmel observed that as urbanism overloads a person, it forces them to distance themselves from the environment \citep{simmel03}, which also may be conducive to differentiation, individualism, and then perhaps freedom. On the other hand, the tendency to hunker down, withdraw, and adopt a blase attitude in cities \citep{simmel03} are not agentic.  
%
% Likewise, there would be more Gemeinschaft in ruralities--more collective and
% social that may restrict personal freedom; and more Gesellschaft in cities that
% may promote individualism (but also hostility and aggression). 
% 
Cities provide freedom from traditional rural social constraints--Gesellschaft versus Gemeinschaft--deviation from community norms is more likely to be reprimanded and ostracized in rural than urban areas. An urbanite may be more internally driven than rural residents, who are more externally constrained \citep{yamagishi12}. 
%rubia was saying sth like that in one of our papers too: 

Cities provide more opportunities for alternative jobs and social relations, so an urbanite is freer by being less concerned about social exclusion as she has plenty of alternatives \citep{yamagishi12}.\footnote{Also see discussion in \citet{luca2023progressive}.} %just srch for 'free' or 'freedom' 
 %rubia: i added as footnote, but we can go back and add more from the paper discussion and then more that to the body of the paper.
 One mechanism is higher social pressure to conform in rural areas--cities provide anonymity--indeed, concern for moral reputation is greater in rural than urban areas
 \citep[Danielson et al 2023 cited in][]{aokOttCh26}. The higher the population density, the lower the frequency of social comparison, as measured by Google searches\citep[][]{aokOttCh26}.  
 
Conversely, there is also theory suggesting that people may feel less free in cities. City life is stressful, hectic, anxiety-provoking, materialistic, consumerist \citep{aokCityBook15}, and overstimulating \citep{simmel03,lederbogen11}--under such conditions, feeling free can be difficult.
Steve Pile in his colorful writing argues that the history and past of old cities can be a weight on their inhabitants \citep{pile05,pile05B,pile99}, and old cities carry melancholia \citep{pile05B}. The weight of history and melancholia can inhibit how free one feels.

Several studies have examined predictors of freedom using the same WVS measure as our study 
\citep[e.g.,][]{aok_free_from_to,pitlik2016free,nikolaev2016give}, but they do not examine urbanicity. There are also many studies on the effect of urban-rural on subjective wellbeing (SWB)
\citep[e.g.,][]{sorensen24,sorensen14,knies26,aok21}, but differences in freedom across the urban-rural gradient have been examined in only a handful of studies covering very few countries. These studies offer a mixed picture, examining both general perceptions of agency and control over specific aspects of daily life. \citet{lolle2023differences} finds no differences in agency across communities of different sizes in Denmark, consistent with our results. 
\citet{atherton2024rural} find that psychological well-being--which entails agency as a subcomponent--is lower in the most rural places compared to the most urban places in the USA; however, these effects become statistically insignificant once socio-demographic factors are controlled for. 
\citet{kusnali2025rural} find that women in Indonesia have more decision making agency in many (i.e., household expenditures, food expenditures, time allocation to social activities) but not all (e.g., contraceptive use) aspects of daily life in urban compared to rural areas. Indian women have more autonomy in cities than in rural areas, both in personal and household decisions \citep{mishra2014relative}. 
B\r{a}ge et al. find no urban-rural difference in people's
perceived reproductive agency, i.e. how much freedom of choice and control they have over family planning, in Ethiopia, Kenya, and Zimbabwe \citep{baage2024attitudes} and sub-Saharan Africa \citep{baage2026reproductive}.
Taken together, these findings leave open the question of whether urban-rural differences in perceived personal agency follow a broader pattern acrross countries. 
 
\subsection{Perceived Political Efficacy}

%rubia: move this down into section
%We turn next to political efficacy: the belief that one has a say in politics and that government will respond to one's voice. 
%AI: It's a person's belief that she has a say in politics, a perception that goverment will respond to one's voice. 
 
%\citet{lord2023autonomy} is a summary of their themed issue on ``Autonomy and
%independence in urban environments.''  Rather suburban than rural contrasts for
%urban, and more on culture and life cycle than space.
 %
As the world continues to urbanize, rural areas are disappearing 
%, and so there is also  ``the perceived disappearance of rural childhood and the values
%associated with it'' in Czech republic
 \citep{pospvech2024loss}. Rurality in the developed world are increasingly peripheral, loosing power and voice relative to urban centers and being abandoned and disrespected--there is rural discontent, resentment, and shame (see superb accounts in 
 \citet{hansonCityJournalautumn15,fullerNYT17monD}). This would suggest that rural areas in rich countries may have lower political efficacy--the belief that one has a say in politics and that government will respond to one's voice. This may reflect their sense of being abandoned and disrespected, rather than inherently lower freedom or personal agency. 
 %!!!opposite to our iniitial country level hypothesis that "This difference is more pronounced in less developed countries"
 %
 %what makes one a real rural person: in a classical study, Bell (1994) identified
 %localism (a history of life in rural areas), ruralism (an experience with a rural kind of job), countryism (knowledge of country life) and communalism (interest in things that express a ties to the
 %community) as key elements of belonging to a place.
 %Peripherality, in other words, is defined through long distances and poor infrastructures but
 %does not have to be experienced through them.
 %
 %And ruralities became digitized--children do not play outside in natura that
 %much anymore, but instead on screens just like their urban counterparts. And
 %autonomy is lost to screens--kids are dependent or even addicted to them.
 %
 %
 In general, a loss of power and respect does not necessarily mean a loss of personal autonomy; rural residents still have a good deal of agency. 

 As in the Czech Republic, in the U.S. there is rural resentment, e.g., people feel that policy is imposed on them, rather than developed with or for them; and there is racial resentment where rural Whites resent urban minorities and immigrants getting better treatment than them \citep{diamond23}. %p505
 % handouts and skip the line
 %
 %ruralities more place attachment and nature conservation
 Although rural life involves reliance on community and tight connections, self-reliance, autonomy, and freedom also remain core values \citep{diamond23}.


%leonie 
Rural areas are 'left-behind' \citep{ulrich2021people}, 'places that don't matter' \citep{rodriguez2018revenge}%cited like 4k!!
---they are second-tier areas with limited or deteriorating opportunities. 
  %
When people's opportunities are constrained, their freedom to exercise agency is likewise restricted. Exercising agency requires access to real and substantively different opportunities from which individuals can choose in pursuing lives that are consistent with their own goals and values. Limited opportunities may therefore undermine both people's perceived ability to shape their own lives in
accordance with these goals and values, and their perceived ability to shape the conditions in which their lives are embedded, such as institutions, infrastructure, and the built and natural environment.
 %
We therefore expect people living in 'left-behind' rural areas to perceive themselves as having less influence over the conditions in which they are embedded and that shape their lives, and consequently to perceive lower levels
of political efficacy. %and personal agency.
 
%leonie
Only a few studies have investigated urban-rural differences in political efficacy, suggesting that people from rural areas and small towns show lower levels of political efficacy compared to urbanites, \citep[e.g.,][]{rowland2024political} for the U.S. and \citet[e.g.,][]{luukkonen2022urbanisation} for Europe.
Political efficacy is higher in urban areas and lowest in
rural areas in Europe \citep{garcia2024they,luukkonen2022urbanisation}, and the U.S. \citep{rowland2024political}.
There are also related studies. \citet{margalit2025cultural} %looks like may be important study 
for example, examine political agency in relation to rurality,
conceptualizing ``rural resentment'' as the perception among rural citizens that the political system gives people like them little or no say in what the government does. Across 10 European countries, they consistently find that rural resentment is associated with a higher likelihood of voting for populist parties, while in the United States it is associated with a higher likelihood of having voted for Trump.
%
Political efficacy has been shown to be positively linked to political engagement \citep{accetta23}. But no one has investigated whether this relationship is universal across
 different contexts (such as countries with high or low GDP). Self-efficacy and political efficacy have been shown to have opposing effects on the justification of terrorism and political violence: people who perceive themselves as highly self-efficacious and those who feel they have little political efficacy are less likely to justify such actions \citep{julkif2022self}.

Overall, these studies suggest that rural residents may feel less politicall efficacy than urban resident, even when they regain a sense of personal agency or freedom. Wether these urban-rural differences extend across countries with different economic development remains an open question. Hence, there is a literature gap and a need to test the urban-rural personal agency and political efficacy gradient across countries using an appropriate dataset. Our study is the first to report the this urban-rural gradient across countries showing the personal agency and political efficacy gradient for each country separately. We use a suitable and accurate dataset and control for relevant predictors as described next.


\section{Data and model}

We begin by postulating our research hypotheses: 

\textit{ H1. People in cities are freer than people in rural areas. Hence, urbanites report greater perceived personal agency and political efficacy than residents in rural or smaller places.}

\textit{ H2. The difference between perceived personal agency and political efficacy among urban versus rural residents is more pronounced in less developed countries}.

\textit{ H3. If one is satisfied financially, her perceived personal agency and
  political efficacy will be greater in urban than rural. %Toennies
}

To test these hypotheses, we use the World Values Survey (WVS) Wave 7 from  
 \url{https://www.worldvaluessurvey.org}. The WVS is an outstanding free dataset that is comparable, if not better suited for our study, than the $\$30k+$ closed source Gallup dataset. It is helpful to use open source public data amidst the corporatization of research and academia \citep{mills2012corporatization,cox2013corporatization,millsNYT12fa,CatropaNYT20feb8,schmidlinNYT15oct10}. 
 %
WVS is representative of each country. Countries in the sample and sample sizes are shown in the appendix.

Urban-rural is a gradient, not a dichotomy \citep{aok-ls_fisher16}, but here for
simplicity we use 3 collapsed categories\footnote{In some specifications, we
  also look at 4 or more categories, please refer to appendix.}.   
 %
Hence, the urbanization variable is as follows: (1) rural areas  $<20k$;  (2)  $<20k$ town areas $>500k$; and (3) city areas   $>500k$--urban unhappiness starts around several hundred thousand \citep{aok-ls_fisher16}.

In the appendix we do several robustness checks, including dropping countries with fewer than 50 observations per category.
%We dropped countries where there were fewer than 50 observations in cities $>100k$
%And in appendix as a robustness/sensitivity check with fewer than 50 observations in cities $>500k$

\subsection{Personal Agency}

The personal agency (i.e., freedom) variable reads:
``Some people feel they have completely free choice and control over their lives, while other people feel that what they do has no real effect on what happens to them.  Please use this scale where 1 means
'none at all' and 10 means 'a great deal' to indicate how much freedom of choice and control you feel you have over the way your life turns out.'' 
The validity and reliability of this measure seems robust and the variable has been used successfully in the literature \citep[e.g.,][]{verme09,inglehart08,veenhoven00_b,aok_free_from_to}. It is nonetheless possible that the question may be interpreted slightly differently in different languages/cultures, and some caution is advised. This is also one of the reasons why we also do within-country analyses in the next section. 
%rubia: i added the above text to the main text, if you think we don't need it, we can comment it out again
All variables including controls are listed in the appendix and there are also shown distributions of all the variables. 
We are controlling for %age, age$^2$, male,
married, employment type dummies (full/part time, self-employed, retired, housewife, student, unemployed, and other),
financial satisfaction,\footnote{Financial satisfaction is key--not everyone needs the same amount of money to be financially satisfied and feel free. This is measured as ``How satisfied are you with the financial situation of your household?''} health, kids,
and importance of religion.

Per controls we follow \citet{pitlik2016free,nikolaev2016give} except that we use financial satisfaction in most specifications instead of income as it arguably predicts freedom more than income. We found that it actually correlates higher in our sample, and also
 we don't use their country level democracy index, and use PCGDP \footnote{PCGDP is a macro-level contextual variable representing per capita Gross Domestic Product (often adjusted or paired with Purchasing Power Parity)} in the appendix only. 
 For the main models, we decided against controlling for age and and age$^2$, as it is rather urbanicity predicting age composition than the other way round, hence we would risk overcontrol bias. Likewise, we do not include gender as it likely differs less than age across urban-rural, and again rather predicted by than predicting
 urban-rural. Robustness tests controlling for these variables are presented in the Appendix and show that results are robust.     

We were tempted to add trust, crime victimization and fear of crime, as these correlate with urbanicity and may predict freedom. 
%and ineed inclusion of these yields larger coefficients.
 %
 % ``Have you been the victim of a crime during the past year?'' and
 % ``In the last 12 months, how often have you or your family: Felt unsafe from
 % crime in your own home''  
 %
% Yet we only include trust and crime victimization and fear of crime as last
% oversaturated model,
But eventually we do not include them because those are rather caused by urbanicity, not the other way round, and hence may lead to overcontrol bias:
\begin{quote}
What we really need, to identify/estimate X$->$Y, is to control for other
determinants of Y that are also antecedents of X (so, W$->$X). If we include
controls where the relationship goes the other way (X$->$W), we will exacerbate
bias in our estimate of X$->$Y. (The relevant term for that situation is 'overcontrol bias') \citep[][p. 5]{bartram24}.
\end{quote}
Additional robustness test were conducted resulting in oversaturated models; the addition of these controls yield larger coefficients, although overall results are the same. Refer to the Appendix for more detail.

\subsection{Political Efficacy}
 The  political efficacy item reads ``how much would you
 say the political system in your country allows people like you to have a say in what the government does?''
 % This was a 5-point scale with a higher number
 % signaling less efficacy. This question was fielded from 2018 onward, so  % countries whose fieldwork was completed before the question was added (except % Argentina where the question was fielded in 2017) were dropped from this
 % analysis due to missing data.
 Data for this item were unavailable for a few countries: the United States, Germany, South Korea, Greece, Turkey, Thailand, China, Chile, Ecuador, Serbia, Romania, Russia, and
 Tunisia.\footnote{Although Tunisian respondents were asked this question, the data for Tunisia were not curated properly and therefore could not be analyzed in this study. \url{https://onlinelibrary.wiley.com/doi/full/10.1111/ssqu.13120}} 
 
% leonie
We control for education, as studies show that more educated people report greater political efficacy
\citep{emmenegger2015labour,garcia2024they,marx2016unemployed,rowland2024political}. We
also control for employment status. Compared with those in paid employment,
people in the education field report greater political efficacy
\textbf{leonie,what do you mean here???},
whereas those who are permanently sick or disabled report lower political efficacy\citep{luukkonen2022urbanisation}. Unemployment and/or precarious employment are also negatively associated with political efficacy
\citep{emmenegger2015labour,marx2016unemployed}. 

We control for respondents' financial situation, as those who can make ends meet report greater political efficacy\citep{luukkonen2022urbanisation,marx2016unemployed}. Similarly, higher income is positively associated with political efficacy \citep{emmenegger2015labour,rowland2024political}. We also control for health and religiosity: better health is associated with greater political efficacy \citep{garcia2024they,gidengil2024healthy,shore2019health}, whereas more religious people report lower political efficacy\citep{marx2016unemployed,rowland2024political}.
 %
These variables also differ across urban and rural areas. Studies document urban-rural differences in education \citep{van2021urban,zahl2024spatial}, employment  \citep{e26,ananian2024employment}, financial situation
 \citep{kpatcha2026unlocking,meloni2024rural,finnemann2024urban},
 health \citep{mackinnon2023mapping,menashe2022shifting},
 and religiosity \citep{balazka2020mapping,meulemann2020secularization}. Controlling for these variables therefore allow us to examine whether urban-rural differences in political efficacy remain after accounting for these variables. 

\section{Results}
We begin by reporting in table \ref{m0} mean personal agency and political
efficacy by settlement size and level of development. It is important to
underscore that even finding no difference in personal agency level between
urban and rural areas would be worth reporting, given our expectation of an
urban advantage and positive relationship.

\begin{table}[H]\centering\footnotesize
\caption{\label{m0} Means}
\begin{tabular} {lrrrr|rrrr}   \hline 
       & \multicolumn{4}{c|}{\bf personal agency}          &\multicolumn{4}{c}{\bf political efficacy}\\
       & \multicolumn{4}{c}{country development:}          &\multicolumn{4}{c}{country development:}\\ 
place  &overall&low/lower middle&upper middle&high&overall&low/lower middle&upper middle&high\\ \hline
$<$20k &7.20   &7.11            &7.26        &7.24&2.96   &3.03            &3.06        &2.67\\
20-500k&7.27   &7.21            &7.21        &7.36&2.88   &2.95            &2.83        &2.87\\
$>$500k&7.12   &6.96            &7.40        &6.98&2.89   &2.67            &2.86        &2.99\\ \hline
\end{tabular}\end{table}

Overall, we find that freedom (i.e. personal agency) has lowest levels in the
largest places, including in less developed countries, and the difference is
especially large for countries with high development.
 But in upper middle development countries, it has highest levels in
cities. Overall, however, these differences are small and disappear after
controlling for predictors of freedom.

Political efficacy is also overall slightly lower in largest cities. Unlike
 personal agency, the political efficacy pattern differs considerably by a country's
level of development.
In less developed countries, cities have much less political say, while in upper
middle countries less, but in high developed countries $<$20k is 2.67  and $>$500k is 2.99,
almost the opposite of low/lower development countries:$<$20k is 3.03  and $>$500k is 2.67.

And that last difference holds even when controlling for predictors of political efficacy. 
Margins are plotted in figure \ref{m-fps}. While the confidence intervals are
wide, the interaction between $>$500k and high developed countries is
significant versus the base case of $<$20k and low/lower development
countries. The results for personal agency are insignificant.


\begin{figure}[H]
 \includegraphics[height=6.5in]{m-fps.pdf}\centering
\caption{\label{m-fps} m-fps these results are robust when droping countries
  where $>$500k $<$50obs and $<$20k $<$50obs}
\end{figure}

  
\subsection{Interactions with Financial Satisfaction}

The interactions with financial satisfaction results for both personal agency
and political efficacy are visualized in figure \ref{zz}. We focus here on
visual interpretations using marginsplot, but all underlying regressions are in
the appendix.
 %
Personal agency results in panel a are significant. The relevant significant
interactions we report are for the smallest places $<20k$ (against the base case
of $>500k$) with top category(ies). 

In panel a, personal agency increases with size among respondents in the top 3
and also middle 3 financial-satisfaction categories from $<20k$ to $>500k$, but
decreases for the bottom 3 categories. 
% 
In other words, personal agency differences are magnified in the largest cities
across financial satisfaction groups--financially satisfied respondents report a
larger advantage in freedom over financially dissatisfied respondents in
cities. The corresponding interaction results for political efficacy in panel b are insignificant. 


\begin{figure}[H]
 \includegraphics[height=6.5in]{zz.pdf}\centering
\caption{\label{zz} zz these results are robust when droping countries
  where $>$500k $<$50obs; but not where  $<$20k $<$50obs}
\end{figure}


% We also separate these results by country, however we leave
% it for the future research as sample sizes run small for interactions.
%rubia: i suggested moving it here, is this the reason why it wasnt? would it
%perhaps be worth reporting here with disclaimer?
%aok: i dont know, lets polish the paper a bit more and decide later
% We also separate results by region and income group for financial satisfaction and  postpone these results to appendix\footnote{We find that the results for financial satisfaction in panel a are stronger in East Asia/Pacific and Sub-Saharan Africa (due to sharp increase for the top tertile), and show similar pattern in Europe/Central Asia except that the middle tertile for financial satisfaction is stronger than the top. The Latin America/Caribbean and North America results show the opposite pattern where the bottom tertile has the steepest slope, and in the Middle East/North Africa the middle and top tertiles even reverse and become negative. In South Asia, the slopes are parallel---we find no differences there. Taken together, by income group we find that high income is similar to the overall pattern.}

\section{Discussion and Conclusion}

The expectation that cities provide greater freedom is plausible, but in our
pooled sample without interactions, we found no significant differences when
controlling for country dummies.
% When analyzing countries individually, the ``city air makes you free'' hypothesis is only supported in 16 out of 64 countries.
%rubia: i would add the results from countries to the previous section, so readers can refer to this easily. thoughts?
% In the pooled sample without interactions there are no significant urban rural differences when controlling for country dummies. But there are some significant results separately by country. The hypothesis that cities are more free does not find support
% except in 16 countries out of 64.
It is worth highlighting that finding no urban-rural effect for personal agency in most countries is a finding worth reporting--the expectation from theory and from the literature is that the bigger places have more freedom. %And in 6 countries we find a positive effect. 

For political efficacy, however, we find it to be significantly higher in urban compared to rural areas in rich countries as hypothetized in the ``left behind ruralities'' hypothesis.
We find support for Toennies' idea that: ``in the city and therefore, where general conditions characteristic of the Gesellschaft prevail, only the upper strata, the rich and the cultured, are really active and alive'' \citep[p. 227]{tonnies57}.
There are significant interactions in the pooled sample for financial satisfaction. 

  
Concurrent with results from \citet{aok24swls}, urbanites probably have, in some ways, fuller and freer lives, but perhaps
they realize more what they miss out relative to other urbanites, and so they feel relatively less free. It's possible that the availability of opportunities does not necessarily translate into the ability to use them. Costs, stress, unequal resources may limit one's choices even when possibilities are endless.   

Another mechanism could be that people tend to compare themselves with those who are better off\citep{frey02s} and there are more people to compare to in urban areas. The inequality ladder is taller and the pond is bigger--a small fish in a big pond (city) seems to be less free than a big fish in a small pond (rural area). 

In poorer/less economically developed countries, there are more differences in freedom across the rural-urban gradient, mostly less freedom in smaller places\footnote{(But sometimes more, refer to appendix)}--this makes sense--smaller places in developing countries can be very traditional and very poor--freedom is difficult in such places. In richer countries, on the other hand, one can already be free just about anywhere, there is no need for a large city. This is a similar to findings in the literature with respect to happiness \citep{aokCityBook15}--only in the very poorest places of all such as Sub-Saharan Africa, is happiness greater in larger cities as life is unbearable and lacks necessities such as shelter and clean water in rural areas. Everywhere else urbanites tend to be less happy \citep{aok21}.

Overall, our findings suggest that the freedom associated with urban areas depends on which aspect we consider. Cities offer more opportunities and choices, but these do not necessarily translate into greater perceived freedom and control over one's life. Yet, for political efficacy, we find an urban advantage in rich countries, consistent with the ``left behind ruralities'' hypothesis. In rich countries, rural residents report lower political efficacy, even though personal agency generally shows little urban-rural difference. Hence, they may feel that they have control over their own lives while feeling that they have less say in government decisions. Future research could shed light on our findings by investigating whether rural residents' perception of having less political voice reflect differences in institutional response, political representation, or access to decision making.
%rubia: we probably need a better closing paragraph, i can't think anymore. lol. more tomorrow.


%rubia: i would incorporate all that was written under "1. more recent results" to manuscript. what do you think? maybe not all tables/figures, but a significant amount. 
%rubia: also needs a table with variable definitions, I was going to copy/paste Table 2, but that seems to include the variables that were not included in the main analyzes (Age, gender) and is also missing political efficacy
%rubia: i also think we should move figure 44 from appendix to paper, i really like that section, and it would be something reviewers i imagine will like. the whole 2.4.3 section could be a new section in the resuls
%aok: maybe! lets revisit
 
\bibliography{/home/aok/papers/root/tex/ebib.bib,cityAgency.bib,/home/aok/papers/root/old/2021/swbCityWorld/tex/swbCityWorld.bib}

\setcounter{section}{0}
 \part*{\Huge ONLINE APPENDIX (SOM)}
 \textbf{[note: this section will NOT be a part of the final version of
   the manuscript, but will be available online instead as SOM 'Supplementary Online
   Material']} %hence everything below

\tableofcontents
\addtocontents{toc}{\protect\setcounter{tocdepth}{3}} 

\subsection{Descriptive Statistics}

var des variable definitions table here\\

hist histograms here\\

\subsection{Regressions underlying body results}

\subsubsection{Interactions with level of development}

\begin{scriptsize}
\begin{spacing}{.19} 
\begin{verbatim}
 mixed polSay i.t3##i.iii  i.X025R i.emp satFin health  rel_imp  || cc: i.t3  ,mle //final! yes educ; no mar no kids

Performing EM optimization ...

Performing gradient-based optimization: 
Iteration 0:  Log likelihood = -101253.19  
Iteration 1:  Log likelihood = -101253.19  

Computing standard errors ...

Mixed-effects ML regression                         Number of obs    =  64,862
Group variable: cc                                  Number of groups =      48
                                                    Obs per group:
                                                                 min =     726
                                                                 avg = 1,351.3
                                                                 max =   3,997
                                                    Wald chi2(20)    =  890.11
Log likelihood = -101253.19                         Prob > chi2      =  0.0000

--------------------------------------------------------------------------------
        polSay | Coefficient  Std. err.      z    P>|z|     [95% conf. interval]
---------------+----------------------------------------------------------------
            t3 |
      20-500k  |      -0.04       0.04    -0.96    0.34        -0.11        0.04
       gt500k  |      -0.06       0.06    -1.07    0.28        -0.18        0.05
               |
           iii |
            2  |      -0.01       0.13    -0.09    0.93        -0.26        0.24
            3  |      -0.20       0.14    -1.49    0.14        -0.47        0.06
               |
        t3#iii |
    20-500k#2  |      -0.05       0.05    -1.02    0.31        -0.16        0.05
    20-500k#3  |       0.16       0.06     2.73    0.01         0.05        0.28
     gt500k#2  |      -0.07       0.08    -0.87    0.39        -0.22        0.09
     gt500k#3  |       0.26       0.09     2.87    0.00         0.08        0.44
               |
         X025R |
       Middle  |       0.01       0.01     0.50    0.62        -0.02        0.03
        Upper  |       0.05       0.01     3.78    0.00         0.02        0.07
               |
           emp |
    Part time  |       0.00       0.02     0.04    0.97        -0.03        0.03
Self employed  |      -0.02       0.01    -1.58    0.11        -0.05        0.01
      Retired  |       0.01       0.02     0.37    0.71        -0.03        0.04
    Housewife  |      -0.07       0.02    -4.40    0.00        -0.10       -0.04
     Students  |       0.00       0.02     0.19    0.85        -0.04        0.04
   Unemployed  |      -0.06       0.02    -3.09    0.00        -0.10       -0.02
        Other  |      -0.03       0.05    -0.72    0.47        -0.13        0.06
               |
        satFin |       0.03       0.00    17.07    0.00         0.03        0.04
        health |       0.07       0.01    12.53    0.00         0.06        0.08
       rel_imp |       0.07       0.01    12.29    0.00         0.06        0.08
         _cons |       2.27       0.09    23.94    0.00         2.09        2.46
--------------------------------------------------------------------------------

------------------------------------------------------------------------------
  Random-effects parameters  |   Estimate   Std. err.     [95% conf. interval]
-----------------------------+------------------------------------------------
cc: Independent              |
                   var(2.t3) |       0.02       0.00          0.01        0.03
                   var(3.t3) |       0.03       0.01          0.02        0.06
                  var(_cons) |       0.13       0.03          0.09        0.20
-----------------------------+------------------------------------------------
               var(Residual) |       1.32       0.01          1.31        1.34
------------------------------------------------------------------------------
LR test vs. linear model: chi2(3) = 4821.97               Prob > chi2 = 0.0000

Note: LR test is conservative and provided only for reference.

. 
. margins,  at(t3=(1 2 3)  iii=(1 2 3))

Predictive margins                                      Number of obs = 64,862

Expression: Linear prediction, fixed portion, predict()
1._at: t3  = 1
       iii = 1
2._at: t3  = 1
       iii = 2
3._at: t3  = 1
       iii = 3
4._at: t3  = 2
       iii = 1
5._at: t3  = 2
       iii = 2
6._at: t3  = 2
       iii = 3
7._at: t3  = 3
       iii = 1
8._at: t3  = 3
       iii = 2
9._at: t3  = 3
       iii = 3

------------------------------------------------------------------------------
             |            Delta-method
             |     Margin   std. err.      z    P>|z|     [95% conf. interval]
-------------+----------------------------------------------------------------
         _at |
          1  |       2.97       0.09    33.11    0.00         2.79        3.14
          2  |       2.96       0.09    32.86    0.00         2.78        3.13
          3  |       2.77       0.10    26.97    0.00         2.56        2.97
          4  |       2.93       0.10    30.80    0.00         2.75        3.12
          5  |       2.87       0.09    30.20    0.00         2.68        3.05
          6  |       2.89       0.11    26.58    0.00         2.68        3.11
          7  |       2.90       0.11    27.41    0.00         2.70        3.11
          8  |       2.82       0.10    28.10    0.00         2.63        3.02
          9  |       2.96       0.11    26.19    0.00         2.74        3.19
------------------------------------------------------------------------------

. marginsplot,saving(a1,replace)ytitle("political say")scheme(s2mono)title("") text(2.98 1.4 "low and lower middle")text(2.9 1.4 "uper middle")
> text(2.78 1.4 "high")recast(line) legend(off)xtitle("size of a place") //xsc(off)

Variables that uniquely identify margins: t3 iii
file a1.gph saved


. mixed free i.t3##i.iii  mar i.emp satFin health kids rel_imp  || cc:  i.t3 ,mle

Performing EM optimization ...

Performing gradient-based optimization: 
Iteration 0:  Log likelihood = -181550.65  
Iteration 1:  Log likelihood = -181550.65  

Computing standard errors ...

Mixed-effects ML regression                        Number of obs    =   84,775
Group variable: cc                                 Number of groups =       61
                                                   Obs per group:
                                                                min =      849
                                                                avg =  1,389.8
                                                                max =    4,018
                                                   Wald chi2(20)    = 10564.59
Log likelihood = -181550.65                        Prob > chi2      =   0.0000

--------------------------------------------------------------------------------
          free | Coefficient  Std. err.      z    P>|z|     [95% conf. interval]
---------------+----------------------------------------------------------------
            t3 |
      20-500k  |       0.01       0.06     0.17    0.86        -0.10        0.12
       gt500k  |       0.09       0.08     1.04    0.30        -0.08        0.25
               |
           iii |
            2  |       0.12       0.17     0.74    0.46        -0.20        0.45
            3  |      -0.03       0.18    -0.16    0.87        -0.37        0.32
               |
        t3#iii |
    20-500k#2  |       0.03       0.07     0.40    0.69        -0.12        0.18
    20-500k#3  |       0.05       0.08     0.67    0.50        -0.10        0.21
     gt500k#2  |       0.04       0.11     0.36    0.72        -0.17        0.24
     gt500k#3  |      -0.12       0.12    -1.00    0.32        -0.35        0.11
               |
           mar |       0.02       0.02     1.38    0.17        -0.01        0.06
               |
           emp |
    Part time  |      -0.09       0.03    -3.31    0.00        -0.15       -0.04
Self employed  |      -0.01       0.02    -0.44    0.66        -0.06        0.04
      Retired  |      -0.06       0.02    -2.33    0.02        -0.11       -0.01
    Housewife  |      -0.20       0.02    -8.49    0.00        -0.25       -0.16
     Students  |      -0.21       0.03    -6.12    0.00        -0.27       -0.14
   Unemployed  |      -0.11       0.03    -3.71    0.00        -0.17       -0.05
        Other  |      -0.17       0.07    -2.49    0.01        -0.31       -0.04
               |
        satFin |       0.26       0.00    79.88    0.00         0.25        0.26
        health |       0.34       0.01    38.79    0.00         0.33        0.36
          kids |       0.03       0.01     4.58    0.00         0.01        0.04
       rel_imp |       0.03       0.01     3.81    0.00         0.02        0.05
         _cons |       4.16       0.13    30.94    0.00         3.90        4.42
--------------------------------------------------------------------------------

------------------------------------------------------------------------------
  Random-effects parameters  |   Estimate   Std. err.     [95% conf. interval]
-----------------------------+------------------------------------------------
cc: Independent              |
                   var(2.t3) |       0.03       0.01          0.02        0.06
                   var(3.t3) |       0.06       0.02          0.03        0.11
                  var(_cons) |       0.27       0.05          0.19        0.39
-----------------------------+------------------------------------------------
               var(Residual) |       4.22       0.02          4.18        4.26
------------------------------------------------------------------------------
LR test vs. linear model: chi2(3) = 4153.75               Prob > chi2 = 0.0000

Note: LR test is conservative and provided only for reference.

. margins,  at(t3=(1 2 3)  iii=(1 2 3))

Predictive margins                                      Number of obs = 84,775

Expression: Linear prediction, fixed portion, predict()
1._at: t3  = 1
       iii = 1
2._at: t3  = 1
       iii = 2
3._at: t3  = 1
       iii = 3
4._at: t3  = 2
       iii = 1
5._at: t3  = 2
       iii = 2
6._at: t3  = 2
       iii = 3
7._at: t3  = 3
       iii = 1
8._at: t3  = 3
       iii = 2
9._at: t3  = 3
       iii = 3

------------------------------------------------------------------------------
             |            Delta-method
             |     Margin   std. err.      z    P>|z|     [95% conf. interval]
-------------+----------------------------------------------------------------
         _at |
          1  |       7.15       0.13    57.16    0.00         6.90        7.40
          2  |       7.27       0.11    65.26    0.00         7.06        7.49
          3  |       7.12       0.12    57.64    0.00         6.88        7.36
          4  |       7.16       0.13    53.85    0.00         6.90        7.42
          5  |       7.31       0.12    62.48    0.00         7.08        7.54
          6  |       7.19       0.13    55.87    0.00         6.93        7.44
          7  |       7.24       0.15    49.35    0.00         6.95        7.52
          8  |       7.40       0.12    59.72    0.00         7.16        7.64
          9  |       7.09       0.13    52.58    0.00         6.83        7.35
------------------------------------------------------------------------------

. marginsplot,saving(a2,replace)ytitle("freedom")scheme(s2mono)title("") text(7.25 2.7 "low and lower middle")text(7.33 1.4 "uper middle")text(
> 7.1 2.7 "high")recast(line) legend(off)xtitle("size of a place") xsc(off)

Variables that uniquely identify margins: t3 iii
file a2.gph saved

graph combine a2.gph a1.gph, cols(1)  xsize(2) ysize(4) scale(1.1)
gr export m-fps.pdf,replace
\end{verbatim}
\end{spacing}{.19}
\end{scriptsize}


\subsubsection{Interactions with Financial Satisfaction}



\begin{scriptsize}
\begin{spacing}{.19} 
\begin{verbatim}
. mixed free i.t3##i.sf3    mar i.emp health kids rel_imp  || cc:   ,mle

Performing EM optimization ...

Performing gradient-based optimization: 
Iteration 0:  Log likelihood = -184338.27  
Iteration 1:  Log likelihood = -184338.27  

Computing standard errors ...

Mixed-effects ML regression                        Number of obs    =   85,997
Group variable: cc                                 Number of groups =       62
                                                   Obs per group:
                                                                min =      849
                                                                avg =  1,387.0
                                                                max =    4,018
                                                   Wald chi2(19)    = 10247.87
Log likelihood = -184338.27                        Prob > chi2      =   0.0000

--------------------------------------------------------------------------------
          free | Coefficient  Std. err.      z    P>|z|     [95% conf. interval]
---------------+----------------------------------------------------------------
            t3 |
      20-500k  |      -0.10       0.04    -2.80    0.01        -0.17       -0.03
       gt500k  |      -0.06       0.04    -1.56    0.12        -0.14        0.02
               |
           sf3 |
            2  |       0.67       0.03    21.86    0.00         0.61        0.73
            3  |       1.51       0.03    45.85    0.00         1.45        1.58
               |
        t3#sf3 |
    20-500k#2  |       0.20       0.04     4.69    0.00         0.12        0.28
    20-500k#3  |       0.14       0.05     3.08    0.00         0.05        0.23
     gt500k#2  |       0.15       0.05     3.15    0.00         0.06        0.24
     gt500k#3  |       0.16       0.05     3.16    0.00         0.06        0.26
               |
           mar |       0.02       0.02     1.24    0.21        -0.01        0.05
               |
           emp |
    Part time  |      -0.09       0.03    -3.38    0.00        -0.15       -0.04
Self employed  |      -0.01       0.02    -0.50    0.62        -0.06        0.03
      Retired  |      -0.05       0.02    -2.17    0.03        -0.10       -0.01
    Housewife  |      -0.21       0.02    -8.89    0.00        -0.26       -0.16
     Students  |      -0.21       0.03    -6.10    0.00        -0.27       -0.14
   Unemployed  |      -0.14       0.03    -4.83    0.00        -0.20       -0.08
        Other  |      -0.21       0.07    -3.06    0.00        -0.35       -0.08
               |
        health |       0.37       0.01    42.11    0.00         0.35        0.39
          kids |       0.03       0.01     4.45    0.00         0.01        0.04
       rel_imp |       0.04       0.01     4.11    0.00         0.02        0.05
         _cons |       4.89       0.08    59.41    0.00         4.73        5.05
--------------------------------------------------------------------------------

------------------------------------------------------------------------------
  Random-effects parameters  |   Estimate   Std. err.     [95% conf. interval]
-----------------------------+------------------------------------------------
cc: Identity                 |
                  var(_cons) |       0.26       0.05          0.18        0.37
-----------------------------+------------------------------------------------
               var(Residual) |       4.25       0.02          4.21        4.29
------------------------------------------------------------------------------
LR test vs. linear model: chibar2(01) = 4212.60       Prob >= chibar2 = 0.0000

. est sto a                

. margins,  at(t3=(1 2 3)  sf=(1 2 3))

Predictive margins                                      Number of obs = 85,997

Expression: Linear prediction, fixed portion, predict()
1._at: t3  = 1
       sf3 = 1
2._at: t3  = 1
       sf3 = 2
3._at: t3  = 1
       sf3 = 3
4._at: t3  = 2
       sf3 = 1
5._at: t3  = 2
       sf3 = 2
6._at: t3  = 2
       sf3 = 3
7._at: t3  = 3
       sf3 = 1
8._at: t3  = 3
       sf3 = 2
9._at: t3  = 3
       sf3 = 3

------------------------------------------------------------------------------
             |            Delta-method
             |     Margin   std. err.      z    P>|z|     [95% conf. interval]
-------------+----------------------------------------------------------------
         _at |
          1  |       6.39       0.07    92.12    0.00         6.26        6.53
          2  |       7.06       0.07   104.88    0.00         6.93        7.19
          3  |       7.91       0.07   115.69    0.00         7.77        8.04
          4  |       6.29       0.07    90.53    0.00         6.16        6.43
          5  |       7.16       0.07   106.88    0.00         7.03        7.29
          6  |       7.95       0.07   116.96    0.00         7.81        8.08
          7  |       6.33       0.07    87.90    0.00         6.19        6.47
          8  |       7.15       0.07   104.65    0.00         7.01        7.28
          9  |       8.00       0.07   114.23    0.00         7.87        8.14
------------------------------------------------------------------------------

. marginsplot, xdimension(t3) xsize(2) ysize(2) saving(z1,replace) text(8.25 4 "{bf:a}", size(large))legend(off)text(7.87 1.32 "top 3 cat")text
> (7 1.32 "middle 3 cat")text(6.3 1.32 "bottom 4 cat")title("")ytitle("personal agency")scheme(s2mono)xsc(off)recast(line) //xtitle("rural-urba
> n")

Variables that uniquely identify margins: t3 sf3
file z1.gph saved

. mixed polSay i.t3##i.sf3   i.X025R  i.emp health  rel_imp  || cc:    ,mle //final! ditched mar and kids and added educ

Performing EM optimization ...

Performing gradient-based optimization: 
Iteration 0:  Log likelihood = -103869.52  
Iteration 1:  Log likelihood = -103869.52  (backed up)

Computing standard errors ...

Mixed-effects ML regression                         Number of obs    =  66,483
Group variable: cc                                  Number of groups =      50
                                                    Obs per group:
                                                                 min =     416
                                                                 avg = 1,329.7
                                                                 max =   3,997
                                                    Wald chi2(19)    =  862.78
Log likelihood = -103869.52                         Prob > chi2      =  0.0000

--------------------------------------------------------------------------------
        polSay | Coefficient  Std. err.      z    P>|z|     [95% conf. interval]
---------------+----------------------------------------------------------------
            t3 |
      20-500k  |      -0.04       0.02    -1.99    0.05        -0.09       -0.00
       gt500k  |      -0.03       0.03    -1.17    0.24        -0.08        0.02
               |
           sf3 |
            2  |       0.10       0.02     5.22    0.00         0.06        0.14
            3  |       0.19       0.02     9.56    0.00         0.15        0.23
               |
        t3#sf3 |
    20-500k#2  |       0.04       0.03     1.35    0.18        -0.02        0.09
    20-500k#3  |       0.04       0.03     1.50    0.13        -0.01        0.10
     gt500k#2  |       0.02       0.03     0.80    0.42        -0.03        0.08
     gt500k#3  |       0.03       0.03     0.81    0.42        -0.04        0.09
               |
         X025R |
       Middle  |       0.00       0.01     0.03    0.98        -0.02        0.02
        Upper  |       0.04       0.01     3.42    0.00         0.02        0.07
               |
           emp |
    Part time  |      -0.00       0.02    -0.18    0.86        -0.04        0.03
Self employed  |      -0.02       0.01    -1.38    0.17        -0.05        0.01
      Retired  |       0.00       0.02     0.17    0.86        -0.03        0.03
    Housewife  |      -0.07       0.02    -4.36    0.00        -0.10       -0.04
     Students  |       0.00       0.02     0.24    0.81        -0.04        0.05
   Unemployed  |      -0.06       0.02    -3.22    0.00        -0.10       -0.02
        Other  |      -0.02       0.05    -0.49    0.62        -0.11        0.07
               |
        health |       0.08       0.01    13.80    0.00         0.06        0.09
       rel_imp |       0.07       0.01    12.79    0.00         0.06        0.08
         _cons |       2.30       0.06    38.08    0.00         2.18        2.42
--------------------------------------------------------------------------------

------------------------------------------------------------------------------
  Random-effects parameters  |   Estimate   Std. err.     [95% conf. interval]
-----------------------------+------------------------------------------------
cc: Identity                 |
                  var(_cons) |       0.13       0.03          0.09        0.20
-----------------------------+------------------------------------------------
               var(Residual) |       1.33       0.01          1.31        1.34
------------------------------------------------------------------------------
LR test vs. linear model: chibar2(01) = 5171.45       Prob >= chibar2 = 0.0000

. est sto a                

. margins,  at(t3=(1 2 3)  sf=(1 2 3))

Predictive margins                                      Number of obs = 66,483

Expression: Linear prediction, fixed portion, predict()
1._at: t3  = 1
       sf3 = 1
2._at: t3  = 1
       sf3 = 2
3._at: t3  = 1
       sf3 = 3
4._at: t3  = 2
       sf3 = 1
5._at: t3  = 2
       sf3 = 2
6._at: t3  = 2
       sf3 = 3
7._at: t3  = 3
       sf3 = 1
8._at: t3  = 3
       sf3 = 2
9._at: t3  = 3
       sf3 = 3

------------------------------------------------------------------------------
             |            Delta-method
             |     Margin   std. err.      z    P>|z|     [95% conf. interval]
-------------+----------------------------------------------------------------
         _at |
          1  |       2.80       0.05    52.26    0.00         2.70        2.91
          2  |       2.90       0.05    55.09    0.00         2.80        3.01
          3  |       3.00       0.05    56.44    0.00         2.89        3.10
          4  |       2.76       0.05    51.09    0.00         2.65        2.87
          5  |       2.89       0.05    54.94    0.00         2.79        3.00
          6  |       3.00       0.05    56.43    0.00         2.89        3.10
          7  |       2.77       0.06    50.20    0.00         2.67        2.88
          8  |       2.90       0.05    54.34    0.00         2.79        3.00
          9  |       2.99       0.05    55.34    0.00         2.89        3.10
------------------------------------------------------------------------------

. marginsplot ,xdimension(t3) xsize(2) ysize(2) saving(z2,replace) text(3.12 4 "{bf:b}", size(large))legend(off)text(2.98 1.42 "top 3 cat")text
> (2.88 1.42 "middle 3 cat")text(2.77 1.42 "bottom 4 cat")title("")ytitle("political efficacy")xtitle("rural-urban")scheme(s2mono)recast(line)

Variables that uniquely identify margins: t3 sf3
file z2.gph saved

. graph combine z1.gph z2.gph, cols(1) scheme(s2mono) xsize(2) ysize(4) scale(1)

. gr export zz.pdf,replace
file zz.pdf saved as PDF format
\end{verbatim}
\end{spacing}{.19}
\end{scriptsize}



\end{spacing}
\end{document}
